% Part: sets-functions-relations
% Chapter: relations
% Section: operations

\documentclass[../../../include/open-logic-section]{subfiles}

\begin{document}

\olfileid{sfr}{rel}{ops}
\olsection{Bewerkingen op relaties}

Het is vaak nuttig relaties te wijzigen of te combineren. In
\olref[sfr][rel][ord]{prop:stricttopartial} hebben we de \emph{vereniging}
van relaties beschouwd. Dit is eenvoudig de vereniging van twee relaties die
we opvatten als verzamelingen van paren. Evenzo hebben we in
\olref[sfr][rel][ord]{prop:partialtostrict} het relatieve verschil van relaties
beschouwd. Hieronder volgen enkele andere bewerkingen op relaties.

\begin{defn}\ollabel{relationoperations} 
Laat $R$ en $S$ relaties zijn en $A$ een willekeurige verzameling.

De \emph{inverse} van $R$ is $R^{-1} = \Setabs{\tuple{y, x}}{\tuple{x,
    y} \in R}$.

De \emph{samenstelling} van $R$ en $S$ (eerst de eerste relatie, daarna de
tweede) is $(R \mid S) =
\{\tuple{x, z} : \exists y(Rxy \land Syz)\}$.

De \emph{restrictie} van $R$ tot $A$ is $\funrestrictionto{R}{A}= R
\cap A^2$.

Het door $R$ bepaalde \emph{beeld} van $A$ is $\funimage{R}{A} = \{y :
(\exists x \in A)Rxy\}$.
\end{defn}

\begin{ex}
Laat $S \subseteq \Int^2$ de opvolgerrelatie op~$\Int$ zijn, dus
$S = \Setabs{\tuple{x, y} \in \Int^2}{x + 1 = y}$. Dan geldt $Sxy$ precies
dan wanneer $x + 1 = y$.

$S^{-1}$ is de voorgangerrelatie op $\Int$, dus
$\Setabs{\tuple{x,y}\in\Int^2}{x -1 =y}$.

$S\mid S$ is
$ \Setabs{\tuple{x,y}\in\Int^2}{x + 2 =y}$

$\funrestrictionto{S}{\Nat}$ is de opvolgerrelatie op~$\Nat$.

$\funimage{S}{\{1,2,3\}}$ is $\{2, 3, 4\}$.
\end{ex}

\begin{defn}[Transitieve afsluiting]Laat $R \subseteq A^2$ een binaire relatie zijn.
	
De \emph{transitieve afsluiting} van~$R$ is $R^+ = \bigcup_{0 < n \in
\Nat} R^n$, waarbij we recursief $R^1 = R$ en $R^{n+1} = R^n
\mid R$ definiëren.

De \emph{reflexief-transitieve afsluiting} van $R$ is $R^* = R^+ \cup
\Id{A}$.
\end{defn}

\begin{ex}
Neem de opvolgerrelatie $S \subseteq \Int^2$. Dan geldt $S^2xy$ precies als
$x + 2 = y$, $S^3xy$ precies als $x + 3 = y$, enzovoort. Dus geldt $S^+xy$
precies als $x + n = y$ voor een zekere $n \geq 1$. Anders gezegd: $S^+xy$
geldt precies als $x < y$, en $S^*xy$ precies als $x \le y$.
\end{ex}

\begin{prob}
Toon aan dat de transitieve afsluiting van $R$ inderdaad transitief is.
\end{prob}
\end{document}
